\section{Tổng quan tài liệu}

\subsection{Nhận diện khuôn mặt dựa trên học biểu diễn sâu}

Các phương pháp nhận diện khuôn mặt hiện đại hoạt động theo mô hình hai giai đoạn: (i) sử dụng mạng nơ-ron sâu để ánh xạ ảnh khuôn mặt vào không gian embedding có chiều thấp, và (ii) so sánh các vector embedding thông qua độ đo cosine similarity để đưa ra quyết định nhận diện. Chất lượng của không gian embedding quyết định trực tiếp hiệu năng hệ thống.

ArcFace~\cite{deng2019arcface} là một trong những phương pháp tiêu biểu nhất cho giai đoạn (i), sử dụng hàm mất mát \textit{additive angular margin} để tối đa hóa khoảng cách góc giữa các lớp danh tính trong không gian embedding. Cùng hướng tiếp cận, CosFace~\cite{wang2018cosface} áp dụng \textit{additive cosine margin} với mục tiêu tương tự. Cả hai phương pháp đều đạt kết quả xuất sắc trên các benchmark chuẩn (LFW > 99,8\%~\cite{huang2008labeled}), tuy nhiên đều được tối ưu cho ảnh đầu vào chất lượng cao và \textbf{không có cơ chế thích nghi tại thời điểm suy luận} khi chất lượng ảnh suy giảm do yếu tố môi trường.

\subsection{Nhận diện khuôn mặt nhận biết chất lượng}

Nhận thấy hạn chế trên, một số nghiên cứu đã tích hợp thông tin chất lượng ảnh vào quá trình nhận diện theo các hướng tiếp cận khác nhau.

\textbf{Thích nghi trong quá trình huấn luyện.} AdaFace~\cite{kim2022adaface} đề xuất sử dụng feature norm $\|\mathbf{z}\|_2$ làm proxy cho chất lượng ảnh, từ đó xây dựng \textit{adaptive margin} trong hàm mất mát: ảnh chất lượng thấp được gán margin nhỏ hơn để giảm tác động tiêu cực lên quá trình tối ưu. Mặc dù hiệu quả trên benchmark hỗn hợp, AdaFace yêu cầu huấn luyện lại (\textit{retraining}) khi gặp điều kiện triển khai mới, điều không khả thi trong bài toán thực tế khi điều kiện ánh sáng thay đổi liên tục theo địa điểm và ca làm việc.

\textbf{Ước tính chất lượng không giám sát.} SER-FIQ~\cite{terhorst2020serfiq} đề xuất ước tính chất lượng ảnh dựa trên độ ổn định của embedding khi sử dụng các \textit{random subnetwork} khác nhau của cùng một mô hình, không yêu cầu nhãn chất lượng (\textit{quality label}). Ảnh chất lượng cao cho embedding ổn định (phương sai thấp), trong khi ảnh chất lượng thấp cho embedding dao động mạnh. Tuy nhiên, SER-FIQ chỉ dừng ở việc tính điểm chất lượng (\textit{quality score}) mà \textbf{không liên kết trực tiếp vào cơ chế ra quyết định ngưỡng} tại thời điểm suy luận.

\subsection{Ngưỡng quyết định thích nghi tại thời điểm suy luận}

Hướng tiếp cận thứ ba tác động trực tiếp vào giai đoạn ra quyết định (ii) bằng cách điều chỉnh ngưỡng thay vì thay đổi mô hình embedding.

Chou et al.~\cite{chou2018adaptive} đề xuất ngưỡng thích nghi theo từng danh tính (\textit{identity-specific threshold}): mỗi người trong gallery được gán một ngưỡng riêng dựa trên phân bố similarity của các mẫu đăng ký. Phương pháp này cải thiện accuracy so với ngưỡng cố định, tuy nhiên có hai hạn chế quan trọng: (i) không xét đến tín hiệu điều kiện môi trường như độ sáng hay độ nhiễu của ảnh truy vấn (\textit{probe image}), và (ii) chi phí tính toán ngưỡng theo từng danh tính tăng tuyến tính với kích thước gallery, gây khó khăn cho triển khai thời gian thực trên thiết bị biên.

\subsection{Khoảng trống nghiên cứu}

Bảng~\ref{tab:related} tổng hợp bốn tiêu chí then chốt mà một hệ thống nhận diện khuôn mặt thích nghi trong triển khai thực tế cần đáp ứng đồng thời: (i) thích nghi tại thời điểm suy luận, (ii) sử dụng tín hiệu môi trường, (iii) không yêu cầu huấn luyện lại, và (iv) khả thi trên thiết bị biên.

\begin{table}[h]
\centering
\caption{So sánh các hướng tiếp cận liên quan theo bốn tiêu chí triển khai thực tế}
\label{tab:related}
\begin{tabular}{lcccc}
\toprule
\textbf{Phương pháp} & \textbf{Inference-time} & \textbf{Env. signal} & \textbf{No retrain} & \textbf{Edge} \\
\midrule
ArcFace~\cite{deng2019arcface}       & \textcolor{badred}{\texttimes}    & \textcolor{badred}{\texttimes}    & \textcolor{badred}{\texttimes}    & \textcolor{goodgreen}{\checkmark} \\
AdaFace~\cite{kim2022adaface}        & \textcolor{badred}{\texttimes}    & \textcolor{badred}{\texttimes}    & \textcolor{badred}{\texttimes}    & \textcolor{badred}{\texttimes}    \\
SER-FIQ~\cite{terhorst2020serfiq}    & \textcolor{goodgreen}{\checkmark} & \textcolor{badred}{\texttimes}    & \textcolor{goodgreen}{\checkmark} & \textcolor{badred}{\texttimes}    \\
Chou et al.~\cite{chou2018adaptive}  & \textcolor{goodgreen}{\checkmark} & \textcolor{badred}{\texttimes}    & \textcolor{goodgreen}{\checkmark} & \textcolor{badred}{\texttimes}    \\
\midrule
\textbf{Nghiên cứu này}              & \textcolor{goodgreen}{\checkmark} & \textcolor{goodgreen}{\checkmark} & \textcolor{goodgreen}{\checkmark} & \textcolor{goodgreen}{\checkmark} \\
\bottomrule
\end{tabular}
\end{table}

Như bảng cho thấy, chưa có nghiên cứu nào kết hợp đồng thời cả bốn tiêu chí: sử dụng tín hiệu môi trường $(L, N)$ để điều chỉnh ngưỡng quyết định tại thời điểm suy luận, không yêu cầu huấn luyện lại mô hình embedding, và triển khai được trên thiết bị biên. Nghiên cứu này nhắm đến việc lấp đầy khoảng trống đó.

Ngoài ra, cần lưu ý rằng ArcFace trong bảng được đánh giá với vai trò là một hệ thống hoàn chỉnh sử dụng ngưỡng cố định (cách triển khai phổ biến nhất), chứ không chỉ với vai trò backbone embedding. Với cách đánh giá này, ArcFace đáp ứng tiêu chí Edge (mô hình \textit{buffalo\_sc} đủ nhẹ cho thiết bị biên) nhưng không thích nghi được theo điều kiện môi trường.